\documentclass[11pt,a4paper]{article}
\usepackage[T1]{fontenc}
\usepackage{lmodern}
\usepackage[top=22mm,bottom=23mm,left=27mm,right=27mm]{geometry}
\usepackage{amsmath,amsthm}
\usepackage[hidelinks]{hyperref}

\hypersetup{
  pdftitle={D2: Existence of a Scapegoat After a Deep Insertion},
  pdfauthor={Parth Pawar}
}
\setlength{\parindent}{0pt}
\setlength{\parskip}{7pt}
\linespread{1.08}
\setlength{\emergencystretch}{1em}
\newcommand{\limitheight}[1]{\left\lfloor\log_{1/\alpha}(#1)\right\rfloor}

\begin{document}

{\centering
  {\large\bfseries D2\par}
  \vspace{5pt}
  {\LARGE\bfseries Existence of a Scapegoat\\[2pt] After a Deep Insertion\par}
  \vspace{9pt}
  {\large Parth Pawar\par}
  \vspace{2pt}
  Structure B Owner \textbar{} CS 253 / MA 253\\
  Search Trees Under Long-Running Churn
  \par\vspace{8pt}
  \rule{35mm}{0.35pt}
\par}
\vspace{6pt}

\section{Setup}

Fix \(\tfrac12<\alpha<1\). Let \(s(v)\) be the number of nodes in the
subtree rooted at \(v\), including \(v\), and let an empty child have
size zero. A node \(v\) is \(\alpha\)-weight-balanced if each child \(c\)
satisfies \(s(c)\leq\alpha s(v)\). Depth counts edges from the root,
which has depth zero. Let \(n\) be the number of nodes immediately after
a successful insertion and before any rebuild.

\textbf{Claim.} If the inserted node has depth \(d>\log_{1/\alpha} n\),
then some ancestor on its path is not \(\alpha\)-weight-balanced.

\section{Proof}

Write the insertion path as \(v_0,v_1,\ldots,v_d\), where \(v_0\)
is the root and \(v_d\) is the new leaf. Thus \(s(v_0)=n\) and
\(s(v_d)=1\). Suppose every edge on this path satisfies
\(s(v_{i+1})\leq\alpha s(v_i)\). Applying the inequality along the path
gives
\[
  1=s(v_d)\leq\alpha s(v_{d-1})
    \leq\cdots\leq\alpha^d s(v_0)=\alpha^d n.
\]
Rearranging gives \((1/\alpha)^d\leq n\). Since \(1/\alpha>1\), taking
the logarithm to base \(1/\alpha\) gives \(d\leq\log_{1/\alpha}n\),
contradicting the insertion depth. Therefore at least one edge satisfies
\[
  s(v_{i+1})>\alpha s(v_i).
\]
The child \(v_{i+1}\) violates the weight bound at its parent \(v_i\).
Hence \(v_i\) is an unbalanced ancestor and can be selected as the scapegoat.
\qed

\section{Applying the test}

Depth is an integer, so the trigger is equivalent to
\(d>\limitheight{n}\). Starting at the new leaf, the implementation
walks upward and chooses the first parent whose child on the path
violates the weight bound. For our choice \(\alpha=2/3\), the test is
\(3s(\mathrm{child})>2s(\mathrm{parent})\). Equality is allowed.

\textbf{Example.} Insert \(1,2,3,4,5\) in order. The new leaf has depth
\(4>\log_{3/2}5\). On the way up, the child of size \(2\) under the
parent of size \(3\) meets the bound. The next edge violates it because
\(3(3)>2(4)\). The scapegoat is key \(2\), whose subtree has four nodes.

\newpage

\section{Height bound used by the validator}

Let \(M\) be the largest size since the last whole-tree rebuild, and let
\(h\) be the largest node depth. For a nonempty tree, the validator checks
\[
  h\leq\limitheight{M}.
\]
The insertion trigger uses the live size \(n\). The validator uses \(M\)
to account for deletions.

\subsection*{Height after a perfect rebuild}

A perfectly balanced subtree of size \(k\geq1\) has the minimum
possible height \(\lfloor\log_2 k\rfloor\). The lower bound follows from
the maximum number of nodes that fit in a given height. For the upper
bound, each child has at most \(\lfloor k/2\rfloor\) nodes, so induction
gives the stated height.

\subsection*{Insertion}
Before any rebuild, the high-water mark can only increase. If the
insertion creates a node deeper than \(\limitheight{M}\), it is the
unique deepest node, and its depth exceeds the bound by exactly one.
Since \(n\leq M\), the insertion triggers the search for a scapegoat.

Let \(r\) be the height of that scapegoat's subtree just before the
rebuild and \(k\) its size. There is only one node at local depth \(r\),
so \(k\leq 2^r\). If \(k=2^r\), every shallower level is full, and the
two child subtrees at the scapegoat differ in size by one. Each child
then has at most \(k/2\) nodes, contradicting the weight violation.

Therefore \(k<2^r\), and a perfect rebuild reduces the height of this
subtree by at least one. All nodes outside it already met the bound.
If an insertion triggers a rebuild without exceeding the high-water
height bound, the minimum-height rebuild preserves the bound.

\subsection*{Deletion and whole-tree rebuilds}
Deletion does not increase height. When the size falls below
\(\alpha M\), the implementation rebuilds the whole nonempty tree
and resets \(M=n\). A whole-tree rebuild caused by insertion also resets
\(M=n\). In either case the resulting height is
\(\lfloor\log_2 n\rfloor\leq\limitheight{n}\).
The one-node tree starts at height zero, so these cases establish the
bound for every nonempty state.

After each update, \(n\geq\alpha M\)
when the tree is nonempty, so
\[
  h\leq\limitheight{M}\leq\limitheight{n}+1.
\]
The argument applies to both rebuild methods because both produce a
perfectly balanced subtree.

\section*{Reference}

I. Galperin and R. L. Rivest.
\href{https://people.csail.mit.edu/rivest/pubs/GR93.pdf}{\emph{Scapegoat Trees}}.
Proceedings of the Fourth Annual ACM--SIAM Symposium on Discrete
Algorithms, 1993, pp.~165--174. See Lemma~5.1 for the existence argument
and Lemma~5.6 for the height reduction after rebuilding.

\end{document}
